\section*{Bab \ref{ch:b3:perturbation-theory} --- Teori Gangguan}

\begin{solution}{exo:b3:perturbation-theory:1}
(a) Tiap $|\varphi_n|^2$ setangkup terhadap $a/2$, sehingga $\langle
x\rangle = a/2$ dan $\Delta E_n = \lambda/2$ untuk semua $n$. (b)
Simetrinya, bukan dinamikanya. (c) Pergeseran yang sama saling
meniadakan dalam tiap selisihnya: spektroskopi, yang mengukur jarak
arasnya, tidak melihat apa pun pada orde pertama. (d) Ketika $\lambda$
menyaingi jarak arasnya, percampuran keadaan yang diabaikannya (orde
kedua) menjadi berperan.
\end{solution}

\begin{solution}{exo:b3:perturbation-theory:2}
$\bra{n\pm1}\hat x\ket n = \sqrt{\hbar/2m\omega}\,\sqrt{n + \tfrac12
\pm \tfrac12}$; kedua suku orde keduanya memberikan
\[
\frac{q^2\mathcal E^2\hbar}{2m\omega}\Big[\frac{n+1}{-\hbar\omega}
+ \frac{n}{\hbar\omega}\Big] = -\frac{q^2\mathcal E^2}{2m\omega^2} ,
\]
yang tidak bergantung pada $n$ --- yakni jawaban eksaknya. Keadaannya
\emph{memang} terkoreksi pada orde pertama (campuran $n \pm 1$
menggeser fungsi gelombangnya menyamping); hanya bagian orde pertama
energinya yang lenyap, karena paritasnya.
\end{solution}

\begin{solution}{exo:b3:perturbation-theory:3}
(a) $E_\pm = \bar E \pm \sqrt{(\Delta/2)^2 + v^2}$ dengan $\Delta =
E_2 - E_1$. (b) $E_- \approx E_1 - v^2/\Delta$, $E_+ \approx E_2 +
v^2/\Delta$: yakni rumus orde keduanya, suku demi suku. (c) Celahnya
$2\sqrt{(\Delta/2)^2 + v^2} \ge \Delta$, dengan kesamaannya hanya pada $v =
0$. (d) Penjabarannya membesar tanpa batas; jawaban eksaknya
memberikan $\pm|v|$ --- yang persis diresepkan oleh pendiagonalan
$\hat V$ dalam blok berdegenerasinya
(\cref{thm:b3:perturbation-theory:degenerate}).
\end{solution}

\begin{solution}{exo:b3:perturbation-theory:4}
(a) Penjabarannya memberi $\ket n$ koefisien $1$; penormalan sampai
orde pertama memaksa koreksinya masuk ke komplemen ortogonalnya. (b)
Dari proyeksinya pada $\bra m$; percampuran yang kecil menuntut
$|V_{mn}| \ll |E_n - E_m|$. (c) $(0.1)^2 = 1\%$. (d)
$(0.1/0.01)^2 = 100$: tak masuk akal --- diagonalkanlah blok dua
arasnya secara eksak.
\end{solution}

\begin{solution}{exo:b3:perturbation-theory:5}
(a) $\hat x^3$ ganjil: rata-ratanya dalam keadaan eigen paritas mana pun
lenyap. (b) Menuliskan $\hat x^4 \propto (\hat a + \hat
a^\dagger)^4$ lalu mempertahankan keenam suku dengan penaikan dan
penurunan yang sama banyak memberikan $\langle(\hat a + \hat a^\dagger)^4\rangle_n = 6n^2
+ 6n + 3$: yakni pergeseran yang dinyatakan itu. (c) Dengan $\gamma < 0$
pergeserannya menjadi makin negatif seperti $n^2$: jarak aras
berurutannya $E_{n+1} - E_n$ menurun secara linear dalam $n$. (d) Anak
tangga tangga getaran sungguhan yang saling mendekat --- nada atasnya
sedikit kurang dari kelipatan nada dasarnya, dan adanya batas disosiasi
yang berhingga.
\end{solution}

\begin{solution}{exo:b3:perturbation-theory:6}
(a) Pasangannya berdegenerasi: penyebutnya lenyap. (b) Unsur
diagonalnya $\lambda(a/2)^2$; unsur taksdiagonalnya $\lambda c^2$ dengan $c =
-8a^2/9\pi^2$ --- integralnya terfaktorkan menjadi kedua bagian
satu dimensinya. (c) Pergeserannya $\lambda[(a/2)^2 \pm
(8a^2/9\pi^2)^2/a^2\cdot a^2]$ --- yakni $\lambda a^2(\tfrac14 \pm
64/81\pi^4)$ --- dengan keadaan eigennya $(\ket{1,2} \pm \ket{2,1})/
\sqrt2$. (d) Simetrinya sudah menamai keadaan eigennya lebih dahulu
(yakni gabungan yang genap dan yang ganjil terhadap penukarannya);
gangguannya hanya dapat memilih keduanya, dan memang itulah yang
dilakukannya.
\end{solution}

\begin{solution}{exo:b3:perturbation-theory:7}
(a) $3ea_0\mathcal E = 3 \times \num{5.29e-11} \times \num{2.5e6}
= \qty{0.40}{meV}$ ke tiap arah. (b) $\bra{100}\hat z\ket{100} = 0$
karena paritasnya, dan tak ada pasangan berdegenerasi yang dapat
dihibridkan dengannya. (c)
$\alpha \sim 2e^2a_0^2/E_{\text{I}} \to \alpha/4\pi\varepsilon_0
\sim 2a_0^3$, yakni orde yang tepat di samping nilai eksaknya $4.5\,a_0^3$. (d)
Tetapan dielektriknya: $\varepsilon_{\text{r}}$ tiap kapasitor adalah
atom yang menjawab medan lewat penyerahan orde kedua ini persis.
\end{solution}

\begin{solution}{exo:b3:perturbation-theory:8}
(a) $A = \omega^3|d|^2/3\pi\varepsilon_0\hbar c^3 \approx
\qty{6e8}{s^{-1}}$: $\tau = \qty{1.6}{ns}$. (b) $\omega$ lebih kecil
sebesar $4.8$: $\omega^3$ sebesar $110$; dengan dipolnya yang agak
lebih besar, garis D-nya keluar pada $\sim\qty{16}{ns}$ --- sesuai
pengukurannya. (c)
$1/A = \qty{3.4e14}{s} \approx 11$ juta tahun. (d) $\omega^3$
membentang $10^{20}$ antara ultraungu dan radio: kilatan UU yang cepat,
kesabaran radio yang berskala geologi --- sehingga \omterm{ex:b3:spin-two-level:hyperfine}{garis 21 cm}
diramalkan (van de Hulst, 1944) dari teori lalu dicari di langit,
tempat $10^{66}$ atom mengimbangi kesabarannya.
\end{solution}

\begin{solution}{exo:b3:perturbation-theory:9}
(a) Kedua keadaannya genap dikalikan $\hat z$ yang ganjil: integrannya
ganjil. (b) Integral $\varphi$-nya $\int_0^{2\pi}\eu^{\iu m\varphi}\dd\varphi$
lenyap kecuali bila $m = 0$. (c) Fotonnya membawa satu $\hbar$ dan
paritas ganjil: $l$ harus berubah satu, dan $m$ paling banyak satu. (d)
Dengan tiap pintu satu foton tertutup, $2s$ hidup $\sim\qty{0.12}{s}$ ---
delapan orde melampaui $2p$ --- lalu meluruh lewat pemancaran serentak
\emph{dua} foton, yakni kontinum lemah yang memang teramati dalam nebula
planet.
\end{solution}

\begin{solution}{exo:b3:perturbation-theory:10}
(a) Mengganti penyebutnya dengan yang terkecil, $\tfrac34
E_{\text{I}}$, lalu memakai ketertutupannya: $\alpha/4\pi\varepsilon_0 \le
2\,a_0^2\,(2E_{\text{I}})\,/\,(\tfrac34 E_{\text{I}})\cdot\tfrac12
= \tfrac{16}3\,a_0^3$. (b) Nilai eksaknya $4.5\,a_0^3$ berada di bawah
batasnya, sebagaimana seharusnya. (c) $n - 1 = N\alpha_{\text{SI}}/2\varepsilon_0
\approx \num{e-4}$ berbanding nilai terukur \num{1.3e-4}: kelunakan
orde kedua atomnya \emph{memang} pembiasan gasnya. (d) Satu jumlahan
unsur matriks menetapkan bagaimana sebuah atom menyerah kepada medan,
dan penyerahan kolektif $10^{25}$ atom tiap meter kubik membelokkan
cahayanya: \omterm{thm:b3:perturbation-theory:corrections}{teori gangguan} $\to$ keterpolarisasian $\to$ optika.
\end{solution}

\begin{solution}{exo:b3:perturbation-theory:11}
(a) Dua cabang hiperbola, dengan pendekatan terdekatnya $2|v|$ pada $t = 0$.
(b) Dalam keadaan yang mengikuti kurva bawahnya secara sinambung ---
keadaan ``diabatik'' awalnya sudah bertukar watak: yakni pemindahan
sepenuhnya. (c) Waktu perlintasannya $\sim v/\lambda$ yang panjang
terhadap $\hbar/v$: adiabatik ketika $v^2/\hbar\lambda \gg 1$. (d)
Kubit: sapukanlah medan kendalinya secara lambat melewati perpotongan
terhindarnya dan populasinya akan menunggangi cabang bawahnya dari satu
keadaan basis ke keadaan basis lainnya --- yakni pemindahan yang tahan
terhadap galat pewaktuan, yang dipakai sehari-hari dalam penyiapan
keadaan secara adiabatik.
\end{solution}

\begin{solution}{exo:b3:perturbation-theory:12}
(a) Tiap suku jumlahannya berupa kuadrat dibagi penyebut yang negatif.
(b) Efek antaratom terdepan dua atom dalam keadaan dasar berorde dua
dalam kopling dipol--dipolnya: negatif --- jadi selalu menarik, bagi
pasangan jenis apa pun. (c) $\bra{\psi_0}\hat H\ket{
\psi_0} = E_0 + \bra{\psi_0}\hat V\ket{\psi_0}$ adalah energi sebuah
keadaan coba, sehingga $\ge$ energi dasar yang sebenarnya: orde
pertamanya hanya dapat melampaui. (d) Keadaan $n = 2$ yang terbelah
Stark adalah keadaan tereksitasi; teoremanya berkenaan dengan keadaan
dasar yang sebenarnya, yang pergeseran linearnya lenyap --- jadi tak ada
pertentangan.
\end{solution}

\begin{solution}{pb:b3:perturbation-theory:1}
\textbf{1.} $\hat V = e\mathcal E\hat z$ tidak mempunyai rata-rata orde
pertama (paritas); jumlahan orde keduanya, yang semua penyebutnya
negatif bagi keadaan dasar, memberikan $-\tfrac12\alpha\mathcal E^2$
dengan $\alpha$ adalah keterpolarisasiannya --- negatif karena arasnya
tertolak dari atas.
\textbf{2.} $\alpha_{\text{SI}} = 4\pi\varepsilon_0 \times
\num{2.2e-30} = \qty{2.4e-40}{F.m^2}$: $p = \qty{2e-37}{C.m}
\approx \num{2e-8}\,ea_0$ --- seperseratus juta dipol atom, tiap
molekulnya.
\textbf{3.} $\hbar\omega \approx \qty{2.3}{eV}$ berbanding celah
ultraungunya $\gtrsim\qty{10}{eV}$: penyebut perturbatifnya nyaris tidak
merasakan penggeraknya, sehingga $\alpha$ statisnya sudah memadai.
\textbf{4.} $p_0 = \alpha\mathcal E_0$ rata terhadap frekuensinya, sehingga $P
\propto \omega^4$: yakni hukum Rayleigh.
\textbf{5.} $(700/450)^4 = 5.9$: cahaya biru terhambur enam kali lebih
banyak daripada merah.
\textbf{6.} Matahari memancarkan lebih sedikit ungu daripada biru dan
kepekaan mata terhadap ungu buruk: langit yang terhambur memuncak,
\emph{bagi kita}, pada biru.
\textbf{7.} $[\alpha/4\pi\varepsilon_0] = \unit{m^3}$ dan
$\omega^4/c^4 = \unit{m^{-4}}$: hasilnya sebuah luas. Pada \qty{550}{nm}:
$\sigma \approx \qty{7e-31}{m^2}$.
\textbf{8.} $\ell = 1/N\sigma \approx \qty{58}{km}$: setebal beberapa
atmosfer --- udara \emph{hampir} tembus pandang.
\textbf{9.} Biru ($\ell \approx \qty{26}{km}$): $1 -
\eu^{-8/26} = 26\%$ terhambur; merah ($\ell \approx \qty{150}{km}$):
$5\%$. Langitnya dibangun dari selisih itu.
\textbf{10.} Sepanjang $\sim\qty{240}{km}$ lintasan matahari rendahnya,
birunya bertahan $\eu^{-9} \sim \num{e-4}$ sedangkan merahnya menyimpan
$\sim20\%$: Matahari terbenam merah --- dan merah yang sama, yang
dibiaskan cincin senja Bumi, mewarnai Bulan yang gerhana menjadi
tembaga.
\textbf{11.} Sebuah dipol tidak memancarkan apa pun sepanjang sumbunya
sendiri: pada hamburan $90^\circ$, komponen osilasi molekulnya sepanjang
garis pandang tidak dapat menyumbang, dan cahaya yang bertahan
terpolarisasi tegak lurus terhadap bidang Matahari--molekul--mata.
\textbf{12.} Awan memaksakan hamburan berganda, yang mengacak
geometrinya: kompas polarisasinya mati di bawah langit mendung --- dan
itulah yang menjadikan penggunaannya oleh orang Viking sebuah prestasi.
\textbf{13.} \emph{Kedudukan} molekulnya: udara adalah gas acak, dengan
kerapatannya berfluktuasi pada tiap skala --- materi yang tertata
sempurna (idealnya sebuah kristal) hanya menghamburkan ke dalam berkas
biasnya.
\textbf{14.} Penghambur yang acak dan terpisah pada skala panjang
gelombang berjumlah dengan fase acak: suku silangnya merata-rata habis
dan keterangannya berjumlah --- persekongkolan merusaknya menuntut
keteraturan, dan ketakteraturan membatalkannya.
\textbf{15.} $1 - \eu^{-8/58} \approx 13\%$ sinar matahari memberi makan
kubah birunya --- selaras dengan langit yang ribuan kali lebih redup
daripada cakram matahari namun cukup terang untuk membaca.
\textbf{16.} Di dalam sebutir titiknya riak gelombangnya menjumlahkan
\emph{amplitudo}: dayanya $\propto N^2$, yang luar biasa besar dan
nyaris buta panjang gelombang (titiknya jauh lebih besar daripada
$\lambda$): awan menghamburkan segalanya --- jadi putih.
\textbf{17.} Penghambur di bawah panjang gelombang memilih warnanya
($\omega^4$); penghambur di atas panjang gelombang memantulkan secara
geometris lalu memutihkannya: susu, kabut, cat, salju.
\textbf{18.} Di dekat titik kritisnya fluktuasi kerapatannya sendiri
tumbuh sampai berukuran optis: fluidanya menjadi awannya sendiri ---
yakni opalesensi kritis, mekanisme Rayleigh yang dikuatkan sampai
kelegapan.
\textbf{19.} $n - 1 = \num{2.5e25} \times \num{2.4e-40}/(2 \times
\num{8.85e-12}) \approx \num{3.4e-4}$, berbanding nilai terukur
\num{2.9e-4}: riak gelombang ke depannya, yang terpadu dengan berkasnya,
memperlambat fasenya.
\textbf{20.} Pergeseran Stark sebuah aras; birunya langit; dan
pembiasan (beserta fatamorgana dan lensa udara) cahaya --- satu
$\alpha$, tiga gejala.
\textbf{21.} Serapan ozon mewarnai senja di zenit menjadi biru; aerosol
memutihkan cakrawalanya; cahaya yang terhambur berganda menjaga langit
tetap sedikit bercahaya setelah matahari terbenam.
\textbf{22.} Tanpa udara, tanpa dipol: Matahari menyala dalam langit
yang hitam --- sehingga membuktikan lewat ketiadaan bahwa kubah biru
kita adalah sinar matahari yang terhambur, bukan sifat ``angkasa'' yang
berpijar.
\textbf{23.} Butiran debu Mars besar terhadap $\lambda$: butiran itu
menghamburkan cahaya merah secara berdaya guna ke segala arah (sehingga
mewarnai langitnya) sambil mendifraksikan biru paling kuat \emph{ke
depan} --- sehingga langitnya karamel dan lingkar cahaya di sekitar
Matahari yang terbenam berwarna biru: ukuran butirannya membalikkan
logika butir 18.
\textbf{24.} Hamburan Thomson rata terhadap frekuensinya: tanpa
pemilihan warna, tanpa biru --- tetapi tanggapan yang sekadar sebanding
dengan \emph{kerapatan elektron} justru itulah yang diinginkan
kristalografi dan radiografi dari sebuah penyelidik.
\textbf{25.} Satu tetapan orde kedua, $\qty{2.2e-30}{m^3}$, yang
dikuadratkan lalu dibobot $\omega^4$, menghasilkan panjang hamburan
cahaya hijau $\sim\qty{60}{km}$: cukup tembus pandang untuk siang hari,
cukup menghambur untuk sebuah kubah biru, terpolarisasi bagaikan sebuah
ladang dipol yang digerakkan, dan memerahkan tiap senja tepat pada
waktunya.
\end{solution}
